\documentclass[book.tex]{subfiles}

\begin{document}
\chapter{Deep Electronic Schr\"{o}dinger Equations}
\label{chap:deep_schrodinger}
\noindent\rule{11cm}{0.4pt} \\
\textbf{Prerequisites:} \autoref{chap:molecular_hamiltonian}, \autoref{chap:dft}, \autoref{chap:hartree_fock} \\
\textbf{Difficulty Level:} ** \\
\noindent\rule{11cm}{0.4pt}
\vspace{5mm}



The electronic Schrödinger equation can only be solved analytically for the hydrogen atom, and the numerically exact full configuration-interaction method is exponentially expensive in the number of electrons. Quantum Monte Carlo methods provide a potential approach that scales well for large molecules, can be parallelized and has, accuracy only limited by the flexibility of the wavefunction ansatz used. PauliNet \cite{foulkes2001quantum} \cite{hermann2020deep} is a deep-learning wavefunction ansatz that achieves nearly exact solutions of the electronic Schrödinger equation for molecules with up to 30 electrons. PauliNet has a multireference Hartree–Fock solution built in as a baseline, incorporates the physics of valid wavefunctions and is trained using variational quantum Monte Carlo. 
%PauliNet outperforms previous state-of-the-art variational ansatzes for atoms, diatomic molecules and a strongly correlated linear H$_{10}$, and matches the accuracy of highly specialized quantum chemistry methods on the transition-state energy of cyclobutadiene, while being computationally efficient.

\subsection{Representing Correlated Wavefunctions}
Machine learning methods have had strong success at predicting electronic energies electron densities and molecular orbitals. This approach entirely avoids the solution of the Schr$\mathrm{\ddot{o}}$dinger equation, at the price of requiring datasets of preexisting solutions, obtained by DFT or the coupled-cluster methods.

By contrast, the direct representation of correlated wavefunctions with neural networks and their unsupervised training via the variational principle 
% Let's make this a reference instead?
%first proposed by Carleo and Troyer for discrete spin lattice systems, 
is an ab initio approach that requires no preexisting data and has no fundamental limits to its accuracy. It is motivated by the fact that neural networks are universal function approximators and can in theory provide more efficient means for approximating the exponentially scaling complexity of many-body quantum systems. %Initial attempts on lattice systems were later generalized to bosons in real space, and even electrons in real space, but the latter approach does not use a wavefunction ansatz in the form of a Slater determinant, and perhaps for that reason does not reach the accuracy of the baseline Hartree–Fock (HF) method.

\subsection{Deep Jastrow Factors}

PauliNet, a deep-learning QMC approach, replaces existing ad hoc functional forms used in the standard Jastrow factor and backflow transformation with more powerful deep neural network representations. Besides gain in expressive power, the neural network architecture is  designed to encode the physics of valid wavefunctions and incorporates the multireference HF method as a baseline. 

%These physically motivated choices are essential to obtain a method that not only is highly accurate, but also converges robustly, while maintaining computational efficiency. 
This neural network ansatz can substantially outperform the accuracy of state-of-the-art wavefunction ansatzes using a similar number of determinants. Thanks to the trainable backflow ansatz, high accuracy can be obtained with orders of magnitude fewer determinants compared to traditional QMC methods. The method has the asymptotic scaling of $O(N^4)$.
%and we expect that it will be feasible to apply it to much larger systems than is currently possible with existing high-accuracy methods. We demonstrate this with an accurate calculation of the transition-state energy of the 28-electron cyclobutadiene molecule, which was previously achievable only with highly specialized methods.

\begin{figure}
    \centering
    % \includegraphics[width=0.7\linewidth]{figures/Differentiable Physics/Deep Quantum/Deep Electronic Schrodingers/paulinet_ansatz.png}
    \includegraphics[width=0.7\linewidth]{figures/Differentiable Physics/Deep Quantum/Deep Electronic Schrodingers/paulinet Ansatz.png}
    \caption{A Diagram of the PauliNet Ansatz}
    \label{fig:enter-label}
\end{figure}

The FermiNet architecture follows the same basic idea, but differs in one important aspect. This architecture does not encode any physical knowledge about wavefunctions besides the essential antisymmetry, which is compensated by a much larger number of optimized parameters. This difference likely leads to the higher computational cost per iteration. In addition, the architecture is trained substantially longer and as a consequence reaches higher accuracy for some systems.

\begin{figure}
    \centering
    % \includegraphics[width=\linewidth]{figures/Differentiable Physics/Deep Quantum/Deep Electronic Schrodingers/ferminet.png}
    \includegraphics[width=\linewidth]{figures/Differentiable Physics/Deep Quantum/Deep Electronic Schrodingers/Ferminet Ansatz.png}
    \caption{A Diagram of the FermiNet Ansatz}
    \label{fig:enter-label}
\end{figure}

\subsection{Deep Electronic Wavefunction Ansatz}
At the core of the deep-learning approach to the electronic Schrödinger equation is a wavefunction ansatz which incorporates both the well-established essential physics of electronic wavefunctions—Slater determinants, multideterminant expansion, Jastrow factor, backflow transformation and cusp conditions—as well as deep networks capable of encoding the complex features of the electronic motion in heterogeneous molecular systems. The proposed trial wavefunction, 
\begin{align}
\psi_{\theta}(r_1, \dotsc, r_N)
\end{align}
is of the multideterminant Slater–Jastrow–backflow type, where both the Jastrow factor, $J$, and the backflow, $f$, are represented by deep networks with trainable parameters $\theta$,

\begin{align}
\psi_{\theta}(r) &= \mathrm{e}^{\gamma (r)+J_{\theta}(r)}\sum_{p} c_p \det \left [\tilde{\varphi}_{\theta, \mu_p i}^{\uparrow}(r) \right ]\det \left [\tilde{\varphi}_{\theta,\mu_p i}^{\downarrow}(r)\right ]\\ 
\tilde{\varphi}_{\mu i}(r) &= \varphi_{\mu }(r_i){f}_{\theta,\mu i}(r)
\end{align}

While the expressiveness of PauliNet is contained in the Jastrow factor and backflow deep networks, the physics is encoded by the determinant form; the one-electron molecular orbitals, $\phi_{\mu}$; and the electronic cusps, $\gamma$, in the following way.

Every valid electronic wavefunction must be antisymmetric with respect to the exchange of same-spin electrons,

\begin{align}
\psi(\ldots,r_{i},\ldots,r_{j},\ldots) &=-\psi (\ldots ,r_{j},\ldots,r_{i},\ldots)
\end{align}

This anti-symmetry is enforced by the use of the matrix determinant in the ansatz for $\psi_\theta$. The model is trained using the variational approach to quantum mechanics. Note that this technique allows for training without having to utilize training data from a higher level of quantum theory.

\subsection{Robust Deep Jastrow Factor and Backflow}

As we mentioned before, the Jastrow factor and backflow functions are represented by deep networks. The PauliNet architecture constrains these deep networks to be invariant or equivariant to the operation $P_{ij}$ of exchanging same-spin electrons.

\begin{align}
J(P_{ij}r) &= J(\vec{r}), \quad P_{ij}{f}_{\mu i}(\vec{r})=f_{\mu j}(P_{ij}\vec{r})
\end{align}

The SchNet architecture provides a natural framework for providing this invariance. SchNet as we covered earlier is a graph neural network that performs updates based on the distance to neighboring nodes. Here $\chi$ is a trainable convolution and the $x$ are feature vectors

\begin{align}
{{\bf{x}}}_{i}^{(n+1)}:={{\bf{x}}}_{i}^{(n)}+{\boldsymbol{\chi }}_{{\boldsymbol{\theta }}}^{(n)}\left(\left\{{{\bf{x}}}_{{{j}}}^{({{n}})},\{| {{\bf{r}}}_{j}-{{\bf{r}}}_{k}| \}\right\}\right)
\end{align}

The Jastrow factor and backflow functions are given by learnable transformations $\eta_\theta$ and $\kappa_\theta$ applied on top of the learned feature vectors from the final output layer of SchNet

\begin{align}
J ={\eta}_{\theta}\left(\sum_{i}{\bf{x}}_i^{(L)}\right),\quad {\bf{f}}_{i} = \kappa_{\theta}\left({\bf{x}}_i^{(L)}\right).
\end{align}

%\subsection{Ansatz optimization}


%\subsection{PauliNet extension of SchNet}

%\subsection{FermiNet}


\end{document}